\documentclass[12pt]{article}
\usepackage[utf8]{inputenc}
\usepackage[T1]{fontenc}
\usepackage{hyperref}
\title{Kuuk-Yak Language Resources for K-12 Curriculum Development}
\author{Compiled by Claude Sonnet 5 (claude-sonnet-5) for Andr\'as Kornai}
\date{}
\begin{document}
\maketitle

\section{Introduction}
Kuuk Yak (ISO 639-3: uky) is an extinct Paman language of the Cape York Peninsula, Queensland. Only a few people in Pormpuraaw are said to remember parts of it. It was once thought possibly to be a dialect of Kuuk Thaayorre. Very little is known about it. Linguist Barry Alpher is reported to be collecting the fragments for a lexicon and grammar, but no published dictionary turned up.

\section{Primary Educational Resources}
No school materials for Kuuk Yak were found.

\section{Linguistic and Reference Materials}
\begin{itemize}
\item Wikipedia entry on Kuuk Yak: \url{https://en.wikipedia.org/wiki/Kuuk_Yak_language}
\item The neighbouring Kuuk Thaayorre has far better documentation and may help with comparison. Wikipedia entry: \url{https://en.wikipedia.org/wiki/Kuuk_Thaayorre_language}
\item Omniglot page on Kuuk Thaayorre: \url{https://omniglot.com/writing/kuukthaayorre.htm}
\item A Grammar of Kuuk Thaayorre (Alice Gaby, Mouton): \url{https://mitpressbookstore.mit.edu/book/9783110653304}
\end{itemize}
These Kuuk Thaayorre items are about a different language and are listed only as related background.

\section{Religious and Community Resources}
None found for Kuuk Yak itself.

\section{Digital and Audio Resources}
None found.

\section{Curriculum Development Context}
Kuuk Yak is extinct and sparsely recorded. A heritage curriculum would have to rest on the few surviving fragments. It would probably need to be tied to the Kuuk Thaayorre community in Pormpuraaw, with its involvement. Whether the Kuuk Yak fragments can be accessed is unknown.
\end{document}